\section*{Audit only: the cited first-syzygy case}
% Verification record, excluded from the manuscript as requested.
\paragraph{A direct proof of the cited first-syzygy case.}
This particular specialization of the cited theorem also has a short
monomial proof. Write
$G_{ij}=\partial\operatorname{per}_m/\partial X_{ij}$. Suppose
\[
 \sum_{i,j,k,l}a_{ij}^{kl}X_{kl}G_{ij}=0.
\]
If $k\ne i$ and $l\ne j$, choose a matching in $G_{ij}$ that includes
$X_{kl}$. The resulting monomial has $X_{kl}^2$, row $i$ absent,
row $k$ doubled, column $j$ absent, and column $l$ doubled.
These row and column multidegrees identify the coefficient
$a_{ij}^{kl}$ uniquely. Hence it is zero.
If $k=i$ and $l\ne j$, fix any two distinct rows $i,h$ and choose a
monomial using column $l$ in those two rows, column $j$ in no row,
and every other column once. Its coefficient is
$a_{ij}^{il}+a_{hj}^{hl}$. Thus these coefficients are negatives
for every pair of distinct rows. Three rows and characteristic zero
force all of them to vanish. The case $k\ne i$, $l=j$ is symmetric.
Only coefficients $a_{ij}:=a_{ij}^{ij}$ remain.

The coefficient of a permutation monomial is
$\sum_i a_{i,\sigma(i)}$. Comparing two permutations differing by a
two-column switch gives
$a_{ij}+a_{hl}=a_{il}+a_{hj}$. Thus
$a_{ij}=\alpha_i+\beta_j$. The condition that every permutation sum
be zero is $\sum_i\alpha_i+\sum_j\beta_j=0$. The one-dimensional
redundancy $(\alpha,\beta)\mapsto(\alpha+c,\beta-c)$ leaves dimension
$2m-2$. These are exactly the row and column expansion differences
used above. This proves \eqref{R23:resolution_start} in the needed
degree directly, while retaining the literature attribution.

